% einstein_photoelectric_ML.tex
% "Einstein's Heuristic as Statistical Inference: The Photoelectric Effect
%  as Machine Learning, 1887--1916"
% Thomas J. Sargent, 2026

\documentclass[12pt]{article}

\usepackage{amsmath,amssymb}
\usepackage{booktabs}
\usepackage{array}
\usepackage{multirow}
\usepackage[margin=1.2in]{geometry}
\usepackage{setspace}
\usepackage{microtype}
\usepackage{hyperref}
\usepackage{natbib}
\usepackage{parskip}

\onehalfspacing

\hypersetup{
  colorlinks  = true,
  linkcolor   = blue,
  citecolor   = blue,
  urlcolor    = blue
}

\title{Einstein's Heuristic as Statistical Inference:\\
       The Photoelectric Effect as Machine Learning,
       1887--1916}
\author{Thomas J.\ Sargent\thanks{%
  This paper is one of five that read episodes in the
  history of physics and astronomy as exercises in statistical inference and
  machine learning.  \citet{sargent2025sources} sets out the theme: what is
  now called artificial intelligence names practices that scientists have used
  for centuries.  The other four treat Kepler and Newton, Pauli's exclusion principle, the
  exoplanet discovery of Mayor and Queloz, and the accelerating universe of
  Riess, Schmidt and Perlmutter.}}
\date{March 2026}

\begin{document}
\maketitle

\begin{abstract}
In his 1905 ``miraculous year'' paper on the photoelectric effect, Albert
Einstein proposed that electromagnetic radiation consists of discrete energy
quanta~$h\nu$.  His argument had two components: a physical \emph{model}
(the light-quantum hypothesis) and an implied \emph{measurement equation}
relating the observable stopping potential~$V_0$ linearly to the incident
frequency~$\nu$.  This note examines Einstein's reasoning from a modern
statistical perspective.  Setting Millikan's measurements against the
qualitative findings of Lenard (1902) and the discrepant estimates of Hughes
(1912) and Richardson \& Compton (1912), I carry out the ordinary
least-squares regression that the theory implies:
$V_0 = \alpha + \beta\nu + \varepsilon$.  The slope estimate $\hat\beta$
estimates $h/e$, and hence Planck's constant $h$.  Millikan's five-point
sodium regression gives $\hat\beta = 4.128\times10^{-15}$\,V$\cdot$s, which
with the modern elementary charge implies $h = 6.62\times10^{-34}$\,J$\cdot$s,
$0.17\%$ below the defined value; the larger error in Millikan's own published
$h$ is inherited from his oil-drop charge, not from the photoelectric slope.
A maximum likelihood treatment under Gaussian errors recovers the same point
estimates and shows that OLS attains the Cram\'er-Rao bound.  The tabulated
stopping potentials are reconstructions that lie on a fitted line to within a
fraction of a millivolt, so their residuals cannot furnish a standard error.
I take the uncertainty from Millikan's nine replicated two-point sessions
($0.7\%$).  It agrees with the Cram\'er-Rao bound at his instrumental
precision ($0.8\%$).  That bound also explains his central design choice.  The
standard error on $h/e$ falls with the spread of the frequencies, not with the
number of readings.
Millikan's measured slope also agrees to $0.35\%$ with the slope that
Planck's 1901 blackbody constant and Millikan's own elementary charge
predict, a comparison in which no photoelectric datum appears.
I then read the episode in the language of machine learning.  Einstein's
quantum hypothesis supplied the hypothesis class and the inductive bias.
Lenard's feature-importance finding eliminated intensity as an input.  The
regression is a supervised-learning problem in which five training points
suffice because the physics fixed the functional form.  A deep network given
the same data recovers neither the constant nor its interpretation.  The
regression selects a functional form and not a mechanism: Einstein's equation
follows without quantising the electromagnetic field, and photon-counting
experiments settled that question only in 1977.  Theory reduces the sample
size that data must supply, and it also determines which measurement
discriminates.
\end{abstract}
%
%\tableofcontents

%======================================================================
\section{Introduction}

In December 1900, Max Planck introduced the energy quantum $h\nu$ as ``an act
of desperation'' to fit the blackbody spectrum.  Five years later, Einstein
borrowed that quantum and applied it to a second phenomenon---the photoelectric
effect---in a paper that would earn him the 1921 Nobel Prize.  Einstein posited
a falsifiable linear equation relating the stopping potential of emitted
electrons to the incident frequency.  Its slope is $h/e$.  Millikan confirmed
the equation to better than $1\%$ in 1916.

The 1905 paper was also a statistical argument, though not the one later
readers often ascribe to it.  Einstein's route to the light quantum was
thermodynamic.  He showed that in the Wien region the \emph{entropy} of
radiation depends on volume as the entropy of an ideal gas of $n$ independent
particles does.  He inferred that radiation behaves, for these purposes, as a
collection of independent localised quanta.  He did not in 1905 read the
classical intensity $|E|^2$ as a probability density for locating a quantum.
He took that step in the 1920s, in the guise of the \emph{Gespensterfeld} or
ghost field: a wave that guides probabilities without itself carrying energy.
Max Born made that reading the cornerstone of quantum mechanics in 1926
\citep{Born1926}, crediting Einstein's guiding field as he generalised it from
photons to the Schr\"odinger wave function.  Einstein won the 1921 Nobel Prize
for the photoelectric equation.  Born won the 1954 Prize for the probabilistic
interpretation.  Einstein spent the rest of his life opposing it.

A modern data scientist confronted with a table of (frequency, voltage) pairs
would recognise the photoelectric experiment as a two-variable linear regression
problem.  The regression came after the theoretical insight.  Albert
Einstein in 1905 wrote down the model
\begin{equation}\label{eq:einstein}
  V_0 = \frac{h}{e}\,\nu - \frac{W}{e}
\end{equation}
before anyone had measured a reliable $V_0$--$\nu$ slope.  Equation
\eqref{eq:einstein} is a prediction from what Planck himself called
``simply an act of desperation''\footnote{%
  Planck introduced the energy quantum $h\nu$ in December 1900 to derive
  the blackbody radiation spectrum, later describing his own step in a
  letter to R.\,W.\,Wood (7 October 1931): ``Briefly summarized, what I
  did can be described as simply an act of desperation\ldots\ a theoretical
  interpretation had to be found at any cost, no matter how high.''
  \citep{PlanckWood1931}.  Einstein borrowed the quantum $h\nu$ five
  years later and applied it to photoemission.}: that light
arrives in discrete quanta of energy $h\nu$.  The slope coefficient is
$h/e$, where $h$ is Planck's constant and $e$ the elementary charge.
The intercept is $-W/e$, where $W$ is the metal's work function---the energy
cost of extracting an electron from the bulk.

Eleven years later, Robert Millikan \citep{Millikan1916} provided the
definitive empirical confirmation, using freshly scraped alkali-metal surfaces
in high vacuum and narrow-band filtered spectral lines from a mercury arc.
His data, assembled in Table~\ref{tab:millikan_data} below, allow a textbook
OLS regression that recovers $h$ to within 1\% of its modern value.

The episode provides a microcosm of scientific inference:
\begin{enumerate}
  \item \textbf{Model specification.}  Einstein's quantum hypothesis imposes
    linearity of $V_0$ in $\nu$, a slope $h/e$ that is the same for every
    metal, and a metal-specific intercept.
  \item \textbf{Experimental design as data quality control.}  Contaminated
    surfaces, stray light, and a restricted frequency range biased the
    estimates of Lenard, Hughes, and Richardson \& Compton.  Millikan solved
    these problems one by one.
  \item \textbf{Parameter identification.}  Variation in $\nu$ across widely
    separated spectral lines identifies the slope $h/e$.  The intercept
    identifies the work function $W$ only up to the contact electromotive
    force (EMF) between the photoemitter and the collector, a nuisance
    parameter that shifts the intercept but not the slope.
  \item \textbf{Reluctance to accept the model.}  Millikan confirmed
    Einstein's equation while publicly rejecting the light-quantum
    interpretation.
  \item \textbf{Unintended conceptual consequences.}  The picture that
    motivated the equation---a field that fixes probabilities for discrete
    quanta---became, twenty years later, Born's rule \citep{Born1926}.
\end{enumerate}

Section~\ref{sec:history} sketches the experimental history.
Section~\ref{sec:equation} develops the regression framework.
Section~\ref{sec:data} assembles the historical data.
Section~\ref{sec:regression} presents the OLS estimates and diagnostics.
Section~\ref{sec:mle} develops maximum likelihood estimation of Einstein's
parameters and derives the Cram\'er-Rao bound.
Section~\ref{sec:contact} discusses the contact-EMF identification problem.
Section~\ref{sec:twopoint} analyses Millikan's nine two-point determinations.
Section~\ref{sec:comparison} compares historical estimates of~$h$.
Section~\ref{sec:millikan_reluctance} distinguishes inference from
interpretation.
Section~\ref{sec:ml} reinterprets the episode in the language of machine
learning and AI.
Section~\ref{sec:dnn} contrasts Einstein's regression with what a deep
neural network would learn from the data available at each epoch.
Section~\ref{sec:conclusions} concludes.

Appendix~\ref{app:lamb_scully} shows that Einstein's equation follows without
quantising the electromagnetic field, so Millikan's regression selects a
functional form and not a mechanism.
Appendix~\ref{app:antibunching} describes the photon-counting experiments that
settled the mechanism in 1977.
Appendix~\ref{app:ml_implications} draws the consequence for
Section~\ref{sec:ml}.  Theory reduces the sample size that data must supply,
and theory also determines which measurement discriminates.

The same statistical anatomy governs the discovery of cosmic acceleration by
Riess, Schmidt and Perlmutter in 1998, which a companion paper treats in the
same way.  General relativity replaces the quantum hypothesis in fixing the
functional form.  A combination of the Hubble constant and the supernovae's
absolute magnitude replaces the contact EMF as the nuisance parameter that
shifts the level of the diagram and leaves its shape alone.  And the fit again
settles a functional form without settling the physics behind it.  The
difference is that Millikan's slope is a single number, fully identified once
the nuisance parameter is removed, while the Friedmann model has two shape
parameters that remain tangled with each other.  The photoelectric case is the
one in which identification succeeds completely, and reading it that way makes
the harder case legible.

%======================================================================
\section{Experimental History, 1887--1916}
\label{sec:history}

\subsection{Discovery and Qualitative Characterisation}

Heinrich Hertz observed in 1887 that ultraviolet light reduced the sparking
distance in a gap discharge.  Wilhelm Hallwachs (1888) showed that a zinc
plate illuminated by UV light lost negative charge.  These observations
established that light could eject negatively charged particles from metals.

J.\,J.\ Thomson (1899) identified the ejected particles as electrons by
measuring their charge-to-mass ratio.  The same year Philipp Lenard began
systematic experiments, culminating in his landmark 1902 paper
\citep{Lenard1902}.

\subsection{Lenard's Key Finding (1902)}

Lenard's decisive contribution was establishing two qualitative facts:
\begin{enumerate}
  \item The \emph{maximum kinetic energy} (equivalently the stopping potential
    $V_0$) of the emitted photoelectrons is \textbf{independent of light
    intensity}.  Classical wave theory predicted the opposite: more intense
    waves deliver more energy.
  \item $V_0$ depends on the \textbf{frequency} (``colour'') of the incident
    light, with higher frequency producing higher $V_0$.
\end{enumerate}
Lenard's quantitative data were, by his own account, rough: his metal surfaces
were oxidised, vacuum conditions poor, and no attempt was made to separate
the dependence of $V_0$ on $\nu$ from confounding effects.  He received the
1905 Nobel Prize in Physics for his cathode-ray work.

\subsection{Einstein's 1905 Hypothesis}

In March 1905 \citep{Einstein1905} Einstein proposed that each light quantum
carries energy $h\nu$.  A single quantum is absorbed by a single electron.
The electron escapes if $h\nu > W$; the excess appears as kinetic energy,
so the stopping potential satisfies equation~\eqref{eq:einstein}.  At the
time of writing, Einstein noted:
\begin{quote}
  ``As far as I can see, our conception is not at variance with the properties
  of the photoelectric effect observed by Mr.\ Lenard.''
\end{quote}
There were, however, no quantitative measurements of the $V_0$--$\nu$ slope.
The predicted slope $h/e$ is a precise, falsifiable number.  Since the 2019
redefinition of the SI both $h$ and $e$ are exact, so the target is exact too:
$h/e = 4.135667\times10^{-15}$\,V$\cdot$s.  In 1905 the number was not known
to that accuracy, but it was already fixed---by Planck's blackbody fit five
years earlier, with no freedom to adjust it to photoelectric data.

\subsection{Early Experiments (1907--1914)}

Between 1907 and 1914 three groups attempted quantitative verification.
Table~\ref{tab:comparison} summarises the estimates of $h$ from the two that
published numbers precise enough to enter it.

\paragraph{Ladenburg (1907).}  The first quantitative $V_0$--$\nu$
measurements, but confined to a wavelength range of some $60$\,nm
(2000--2600\,\AA).  Over so narrow a span a linear and a square-root law are
nearly indistinguishable, and Ladenburg's data could not discriminate between
Einstein's hypothesis and its classical rivals.  He reported no determination
of $h$ stable enough to tabulate, and none appears in
Table~\ref{tab:comparison}.

\paragraph{Hughes (1912).}  Arthur Llewelyn Hughes \citep{Hughes1913} used
three UV lines ($\lambda = 1849$, 2257, 2537\,\AA) and roughly a dozen metals.
Two of the three frequencies were used to fit the line; the third served as a
check.  Hughes reported $h$ values ranging from $3.55\times10^{-27}$ to
$5.85\times10^{-27}$\,erg$\cdot$s---a spread of 65\%.  Hughes did not measure
or correct the contact potentials.  The $69$\,nm span of his three lines is
too narrow to pin down a slope.  Millikan's six lines span $293$\,nm.

\paragraph{Richardson \& Compton (1912).}  Owen Richardson and Karl Compton
\citep{RichardsonCompton1912} measured eight metals.  Their $h$ estimates
varied by more than 60\%.  Millikan \citeyearpar{Millikan1916} attributed
the scatter to the same deficiencies as Hughes: narrow frequency range,
uncontrolled contact EMF, and stray short-wavelength scattered light.

\subsection{Millikan's Definitive Measurement (1916)}

Robert Millikan spent a decade developing his ``machine shop in vacuo.''
Inside an evacuated glass bulb, an electromagnetically driven knife freshly
shaved the surface of a cast cylinder of sodium, potassium, or lithium.
The wheel rotated the fresh surface into the light beam without breaking vacuum.
The contact EMF was independently measured by the Kelvin double-balance
method and monitored throughout each run.  Only spectral lines with no
short-wavelength companions in the mercury arc were used, and those lines were
further filtered to remove scattered UV.

The result was the regression in Section~\ref{sec:regression}.

%======================================================================
\section{The Regression Equation}
\label{sec:equation}

\subsection{Derivation}

Einstein's light-quantum hypothesis and the energy-conservation condition
for photoelectron escape give
\begin{equation}\label{eq:model}
  \underbrace{eV_0}_{\text{max KE}}
  = \underbrace{h\nu}_{\text{photon energy}}
  - \underbrace{W}_{\text{work function}},
\end{equation}
where all quantities are in SI units.  Dividing by $e$ gives the true
stopping potential as a function of frequency,
\begin{equation}\label{eq:obs}
  V_0 = \underbrace{\frac{h}{e}}_{\beta}\,\nu
        \;\underbrace{-\;\frac{W}{e}}_{\text{work function}}
        + \underbrace{\varepsilon}_{\text{meas.\ error}}.
\end{equation}
The apparatus reads an applied potential displaced by the contact EMF $v_c$
between photoemitter and collector,
\begin{equation}\label{eq:obs_raw}
  V_0^{\mathrm{obs}} = V_0 - v_c
  = \frac{h}{e}\,\nu - \frac{W}{e} - v_c + \varepsilon .
\end{equation}
Because $v_c$ enters additively it shifts the intercept and leaves the slope
alone.  Writing $\alpha$ for whichever intercept is in force, the estimating
equation is
\begin{equation}\label{eq:reg}
  \boxed{V_0 = \alpha + \beta\,\nu + \varepsilon, \qquad
  \beta = \frac{h}{e}.}
\end{equation}
This is a simple bivariate linear regression with one unknown slope ($\beta$)
and one unknown intercept ($\alpha$).

\medskip
\noindent\textbf{Convention.}  Millikan measured $v_c$ independently by the
Kelvin method and added it back before plotting, and the stopping potentials
tabulated in Section~\ref{sec:data} are the corrected values
$V_0 = V_0^{\mathrm{obs}} + v_c$.  Every regression reported below is run on
those corrected values, so throughout this paper
\begin{equation}\label{eq:alpha_meaning}
  \alpha = -\,W/e ,
\end{equation}
and the work function is recovered from the intercept alone,
$W/e = -\hat\alpha$.  Had the raw readings been used instead, the intercept
would have been $-W/e-v_c$, and separating the two would have required an
independent measurement of $v_c$.  Adding $v_c$ to $-\hat\alpha$ when the data
already contain the correction double-counts the contact EMF.

\subsection{Identification}

\paragraph{The slope.}
Variation in $\nu$ across spectral lines identifies $\beta = h/e$.  The
standard error of $\hat\beta$ is $s/\sqrt{S_{\nu\nu}}$, where
$S_{\nu\nu} = \sum_i(\nu_i - \bar\nu)^2$; spreading the $\nu_i$ apart raises
$S_{\nu\nu}$ and shrinks the standard error.  The early experiments were
confined to UV bands some $60$--$70$\,nm wide and gave unreliable estimates;
Section~\ref{sec:fisher} quantifies the penalty.

\paragraph{The intercept and the contact-EMF problem.}
The intercept is contaminated by the contact EMF $v_c$, which was the dominant
experimental nuisance of the whole literature: it could reach $3$\,V,
comparable to the stopping potentials themselves, and it drifted with surface
age, adsorbed gas, and temperature.  Because it enters additively it is
harmless to the slope so long as it is constant within a run, and fatal to it
otherwise.  Section~\ref{sec:contact} takes up both the identification
problem and Millikan's solution to it.

\paragraph{The universal slope as a cross-metal restriction.}
Because $\beta = h/e$ is a universal physical constant, equation
\eqref{eq:reg} implies that \emph{all} metals share the same slope.  This is
a strong and testable restriction.  Economists will recognise a
cross-equation restriction, in which a single deep parameter appears
identically in the reduced form of every metal.  Separate regressions must
return the same $\hat\beta$.  Millikan reported that his sodium and lithium
slopes did.  Section~\ref{sec:panel} imposes the restriction, tests it, and
explains why his data support it less strongly than the raw agreement
suggests.

%======================================================================
\section{The Historical Data}
\label{sec:data}

\subsection{Millikan's Mercury Arc Lines}

Millikan used isolated lines from a mercury arc whose wavelengths were
independently measured by Walter T.\ Whitney by grating photograph.
Table~\ref{tab:hg_lines} lists the lines and their frequencies (computed
from $\nu = c/\lambda$ with $c = 2.998\times10^{10}$\,cm/s).

\begin{table}[ht]
\centering
\caption{Mercury-arc spectral lines used by \citet{Millikan1916}.
Frequencies computed from $\nu = c/\lambda$,
$c = 2.998\times10^{10}$\,cm/s.}
\label{tab:hg_lines}
\begin{tabular}{rrl}
\toprule
$\lambda$ (\AA) & $\lambda$ (nm) & $\nu$ ($10^{14}$\,Hz) \\
\midrule
5461.0 & 546.10 & 5.490 \\
4339.4 & 433.94 & 6.908 \\
4046.8 & 404.68 & 7.409 \\
3650.2 & 365.02 & 8.213 \\
3125.5 & 312.55 & 9.591 \\
2534.7 & 253.47 & 11.827 \\
\bottomrule
\end{tabular}
\end{table}

\subsection{Millikan's Raw Stopping-Potential Data}

Table~\ref{tab:millikan_raw} reproduces the control readings near the
balance point (zero electrometer deflection) from Millikan's Table~I.
The electrometer sensitivity was 2500\,mm/V; the Weston voltmeter was
certified by the National Bureau of Standards to $\pm0.1\%$.

\begin{table}[ht]
\centering
\caption{Selected readings from Millikan (1916), Table~I.
Sodium surface.  ``V'' is the applied potentiometer voltage;
``mm'' is the electrometer deflection in 30\,s.  For lines
5461--3126\,\AA\ the contact EMF ($v_c = 2.51$\,V) exceeds the
photopotential and the reported voltages are the \emph{assisting}
potential near balance (balance = zero deflection).
For line 2535\,\AA\ voltages are retarding.}
\label{tab:millikan_raw}
\begin{tabular}{rr|rr|rr|rr|rr|rr}
\toprule
\multicolumn{2}{c|}{5461\,\AA} &
\multicolumn{2}{c|}{4339\,\AA} &
\multicolumn{2}{c|}{4047\,\AA} &
\multicolumn{2}{c|}{3650\,\AA} &
\multicolumn{2}{c|}{3126\,\AA} &
\multicolumn{2}{c}{2535\,\AA} \\
V & mm & V & mm & V & mm & V & mm & V & mm & V & mm \\
\midrule
2.257 & 28  & 1.581 & 44  & 1.576 & 82   & 1.157 & 67.5 & 0.581 & 52   & $-0.058$ & 68   \\
2.205 & 14  & 1.629 & 20  & 1.524 & 55   & 1.105 & 36   & 0.529 & 29   & $+0.058$ & 38   \\
2.152 &  7  & 1.576 & 10  & 1.471 & 36   & 1.053 & 19   & 0.477 & 12   & $+0.162$ & 26   \\
2.100 &  3  & 1.524 &  4  & 1.419 & 24   & 1.000 & 11   & 0.424 &  5   & $+0.267$ & 16.5 \\
      &     & 1.367 &  3  & 0.948 &  4   &       &      & 0.372 &  2.5 & $+0.372$ &  8   \\
\bottomrule
\end{tabular}
\end{table}

\noindent The balance points interpolated from these curves (Millikan's
Fig.\ 5) are the stopping potentials used in the regression.

\subsection{Reconstructed Stopping Potentials}

Table~\ref{tab:millikan_data} collects the true stopping potentials
$V_{0,\mathrm{true}} = V_{0,\mathrm{obs}} + v_c$ for sodium and lithium.  For
sodium, the contact EMF was independently measured as $v_c = 2.51$\,V
(stable to $\pm 0.01$\,V over 9 hours).  For lithium (first series),
$v_c = 1.52$\,V.  For sodium at 2535\,\AA\ the paper explicitly states
$V_{0,\mathrm{obs}} = 0.52$\,V, giving $V_{0,\mathrm{true}} = 3.03$\,V,
which lies $0.04$\,V below the value implied by the line fitted to the
other five points.

\paragraph{Provenance of the tabulated values.}
The 2535\,\AA\ sodium entry is Millikan's printed figure.  The other five
sodium entries are reconstructions.  I read them off the straight line
Millikan fitted to his Fig.~5 balance points and rounded them to three
decimals.  His plotted curves cannot be read to better than a few hundredths
of a volt.  The lithium column is a reconstruction too.  Its entries at
4339--3126\,\AA\ are the sodium entries less a constant $0.540$\,V.

Two consequences follow.  Residuals computed from these five points measure
rounding, not measurement scatter, and a standard error built from them
understates the uncertainty by more than an order of magnitude.  The sodium
and lithium slopes agree over the four shared lines by construction.

Sections~\ref{sec:regression} and~\ref{sec:mle} therefore exhibit the
inferential logic that Einstein's equation implies.
Section~\ref{sec:twopoint} supplies the uncertainty estimate.  Millikan's
nine two-point sessions are repeated measurements with fresh surfaces.

\begin{table}[ht]
\centering
\caption{Historical stopping-potential data used in the regression.
True stopping potentials $V_0 = V_{0,\text{obs}} + v_c$ for sodium
(Na, $v_c = 2.51$\,V) and lithium (Li, $v_c = 1.52$\,V).
These values are reconstructed from Millikan (1916) as described
in Section~\ref{sec:data}; only the Na entry at 2535\,\AA\ is a directly
quoted figure.  The Li entries at 4339--3126\,\AA\ are the Na entries less a
constant $0.540$\,V, so the two slopes agree over those lines by construction
(Section~\ref{sec:data}).
The dash for Li at 5461\,\AA\ indicates
that this wavelength is above lithium's photoemission threshold
($\lambda_0 \approx 5260$\,\AA).}
\label{tab:millikan_data}
\medskip
\begin{tabular}{rrrcc}
\toprule
$\lambda$ (\AA) & $\lambda$ (nm) & $\nu$ ($10^{14}$\,Hz)
  & $V_0^{\mathrm{Na}}$ (V)
  & $V_0^{\mathrm{Li}}$ (V) \\
\midrule
5461.0 & 546.10 &  5.490 & 0.454  & ---    \\
4339.4 & 433.94 &  6.908 & 1.039  & 0.499  \\
4046.8 & 404.68 &  7.409 & 1.246  & 0.706  \\
3650.2 & 365.02 &  8.213 & 1.578  & 1.037  \\
3125.5 & 312.55 &  9.591 & 2.147  & 1.606  \\
2534.7 & 253.47 & 11.827 & 3.030  & 2.529  \\
\bottomrule
\end{tabular}
\end{table}

%======================================================================
\section{OLS Regression: Estimating $h/e$}
\label{sec:regression}

\subsection{Regression Setup}

The model is
\[
  V_{0i} = \alpha + \beta\,\nu_i + \varepsilon_i,
  \quad \varepsilon_i \overset{\text{iid}}{\sim}(0,\sigma^2),
  \quad i = 1,\ldots,n.
\]
OLS minimises $\sum_{i=1}^{n}(V_{0i} - \alpha - \beta\nu_i)^2$.
The closed-form estimates are
\begin{align}
  \hat\beta &= \frac{S_{\nu V}}{S_{\nu\nu}}
    = \frac{\sum_i(\nu_i-\bar\nu)(V_{0i}-\bar V_0)}
           {\sum_i(\nu_i-\bar\nu)^2}, \label{eq:betahat}\\
  \hat\alpha &= \bar V_0 - \hat\beta\,\bar\nu. \label{eq:alphahat}
\end{align}

\subsection{Results: Sodium (Five Points, $\lambda = 5461$--$3126$\,\AA)}

Millikan's primary regression used the five most reliable sodium points
(omitting 2535\,\AA\ where a scattered-light correction was required).
I reproduce this regression exactly.

\medskip\noindent\textbf{Data.}
\[
\begin{array}{c|cccccc}
\hline
\nu\;(10^{14}\text{ Hz}) & 5.490 & 6.908 & 7.409 & 8.213 & 9.591 \\
V_0\text{ (V)           }& 0.454 & 1.039 & 1.246 & 1.578 & 2.147 \\
\hline
\end{array}
\]

\medskip\noindent\textbf{Summary statistics.}
\begin{align*}
  n &= 5, \\
  \bar\nu &= \tfrac{1}{5}(5.490 + 6.908 + 7.409 + 8.213 + 9.591)
           = 7.522\times10^{14}\text{ Hz}, \\
  \bar V_0 &= \tfrac{1}{5}(0.454+1.039+1.246+1.578+2.147)
            = 1.293\text{ V}.
\end{align*}

\medskip\noindent\textbf{Sums of squares and cross-products.}

\begin{center}
\begin{tabular}{c|rrrr}
\toprule
$i$ & $\nu_i - \bar\nu$ ($10^{14}$\,Hz) & $V_{0i} - \bar V_0$ (V)
  & $(\nu_i-\bar\nu)^2$ & $(\nu_i-\bar\nu)(V_{0i}-\bar V_0)$ \\
\midrule
1 & $-2.032$ & $-0.839$ & $4.129$ & $1.705$ \\
2 & $-0.614$ & $-0.254$ & $0.377$ & $0.156$ \\
3 & $-0.113$ & $-0.047$ & $0.013$ & $0.005$ \\
4 & $+0.691$ & $+0.285$ & $0.477$ & $0.197$ \\
5 & $+2.069$ & $+0.854$ & $4.281$ & $1.767$ \\
\midrule
$\sum$ & & & $S_{\nu\nu} = 9.277$ & $S_{\nu V} = 3.830$ \\
\bottomrule
\end{tabular}
\end{center}

\noindent(Units: $S_{\nu\nu}$ in $(10^{14}\text{ Hz})^2$;
$S_{\nu V}$ in $10^{14}$\,Hz$\cdot$V.)

\medskip\noindent\textbf{OLS estimates.}
\begin{align}
  \hat\beta &= \frac{S_{\nu V}}{S_{\nu\nu}}
             = \frac{3.830}{9.277}
             = \mathbf{0.41285}\times10^{-14}\text{ V/Hz}
             = \mathbf{4.128\times10^{-15}\text{ V$\cdot$s}},
             \label{eq:betahat_num} \\
  \hat\alpha &= \bar V_0 - \hat\beta\,\bar\nu
             = 1.293 - 4.128\times10^{-15}\times7.522\times10^{14}
             = 1.293 - 3.106 = \mathbf{-1.813\text{ V}}.
             \label{eq:alphahat_num}
\end{align}

\medskip\noindent\textbf{Residuals and goodness of fit.}

\begin{center}
\begin{tabular}{c|rrr}
\toprule
$\lambda$ (\AA) & $V_{0i}$ (V) & $\hat V_{0i}$ (V) & $\hat\varepsilon_i$ (mV) \\
\midrule
5461 & 0.454 & 0.4538 & $+0.18$ \\
4339 & 1.039 & 1.0392 & $-0.23$ \\
4047 & 1.246 & 1.2461 & $-0.07$ \\
3650 & 1.578 & 1.5780 & $+0.01$ \\
3126 & 2.147 & 2.1469 & $+0.11$ \\
\bottomrule
\end{tabular}
\end{center}

\noindent The coefficient of determination is
$R^2 = 1 - \mathrm{RSS}/S_{VV} = 1 - (1.03\times10^{-7})/(1.58) \approx
0.9999999$.

\medskip\noindent\textbf{The residuals measure rounding.}
Every residual is smaller than $0.25$\,mV, forty times below Millikan's
stated reading precision of $\pm10$\,mV.  The reconstruction of
Section~\ref{sec:data} produced them.  Five of these six $V_0$ values were
read off a fitted line and rounded to three decimals.  Taken at face value
they give
\[
  \hat\sigma = \sqrt{\frac{\sum\hat\varepsilon_i^2}{n-2}}
  = \sqrt{\frac{1.03\times10^{-7}}{3}} = 0.18\text{ mV},
  \qquad
  \mathrm{SE}(\hat\beta) = \frac{\hat\sigma}{\sqrt{S_{\nu\nu}}}
  = 6.1\times10^{-20}\text{ V$\cdot$s},
\]
a relative precision of $0.0015\%$.  That is three hundred times finer than
Millikan claimed.

\medskip\noindent\textbf{A plug-in from the instrument.}
I take $\sigma$ from the instrument instead of from the residuals.  Millikan's
balance points were reproducible to about $\sigma = 10$\,mV.  Substituting
that value,
\begin{equation}\label{eq:se_honest}
  \mathrm{SE}(\hat\beta) = \frac{\sigma}{\sqrt{S_{\nu\nu}}}
  = \frac{0.010\text{ V}}{\sqrt{9.277}\times10^{14}\text{ Hz}}
  = 3.3\times10^{-17}\text{ V$\cdot$s},
\end{equation}
a relative uncertainty of $0.8\%$, and $0.5\%$ for the six-line design
(Section~\ref{sec:fisher}).  Millikan claimed ``not more than 0.5\%'' from
graphical fitting.  His nine two-point sessions give $0.7\%$
(Section~\ref{sec:twopoint}).  I use $\sigma = 10$\,mV throughout.

\subsection{Results: Sodium (Six Points, Including 2535\,\AA)}

Including the 2535\,\AA\ point ($\nu = 11.827\times10^{14}$\,Hz,
$V_0 = 3.030$\,V---Millikan's explicitly stated true stopping potential)
widens the design by a further $59$\,nm but moves the slope appreciably:

\begin{align*}
  \bar\nu &= 8.240\times10^{14}\text{ Hz},\quad
  \bar V_0 = 1.582\text{ V}, \\
  S_{\nu\nu} &= 24.72\times(10^{14})^2, \quad
  S_{\nu V} = 10.06\times10^{14}\text{ V$\cdot$Hz}, \\
  \hat\beta &= \frac{10.06}{24.72} = 4.070\times10^{-15}\text{ V$\cdot$s},
\end{align*}
$1.4\%$ below the five-point estimate, with a residual standard deviation of
$11$\,mV---two orders of magnitude larger than in the five-point fit, and now
in line with Millikan's stated reading precision.

The single genuine datum in the table drives the movement.  The 2535\,\AA\
point lies $0.04$\,V below the five-point line (Section~\ref{sec:data}), and
its extreme frequency gives it the highest leverage.  Millikan attributed the
deficit to scattered short-wavelength light in the far UV and excluded the
line from his primary determination.  The observation that contributes most to
$S_{\nu\nu}$ also contributes most to any systematic error in the far UV.
Widening the design and protecting it from stray light pull against each
other.  The $11$\,mV residual standard deviation here, against $0.18$\,mV for
the five-point fit, confirms that the five-point residuals were rounding
artefacts.

\subsection{Results: Lithium (Five Points)}

Repeating for lithium using lines 4339--2535\,\AA\ ($n=5$,
$S_{\nu\nu} = 15.65\times(10^{14})^2$, $S_{\nu V} = 6.457\times10^{14}$):

\begin{align*}
  \hat\beta_{\mathrm{Li}} &= 4.127\times10^{-15}\text{ V$\cdot$s},
  \qquad \hat\alpha_{\mathrm{Li}} = -2.352\text{ V},
\end{align*}
which differs from the sodium slope by $0.04\%$.

The agreement is not an independent confirmation of universality.  Four of
the five lithium entries are the sodium entries shifted by a constant
(Section~\ref{sec:data}), and a constant shift leaves the slope unchanged.
Only the 2535\,\AA\ line carries information absent from the sodium column.
The comparison establishes something weaker.  Millikan reported the two
metals' stopping-potential curves as parallel at his stated precision.

\subsection{Deriving Planck's Constant}

With the OLS slope $\hat\beta = 4.128\times10^{-15}$\,V$\cdot$s (five-point
sodium regression) and the elementary charge $e$ from Millikan's oil-drop
experiment:

\begin{center}
\begin{tabular}{p{5.2cm}p{3.8cm}r}
\toprule
Value of $e$ used & Implied $h = \hat\beta \cdot e$ & Error \\
\midrule
Millikan (1916): $1.592\times10^{-19}$\,C
  & $6.573\times10^{-34}$\,J$\cdot$s & $-0.81\%$ \\
Modern: $e = 1.602177\times10^{-19}$\,C
  & $6.616\times10^{-34}$\,J$\cdot$s & $-0.17\%$ \\
Modern exact: $h = 6.626070\times10^{-34}$\,J$\cdot$s & --- & 0 \\
\bottomrule
\end{tabular}
\end{center}

\noindent The error decomposes cleanly.  Paired with the modern $e$, the
slope alone accounts for $-0.17\%$; Millikan's own $e$, too small by
$0.64\%$, contributes the remainder.  His $h$ was therefore in error
predominantly through the oil-drop constant rather than through the
photoelectric slope.  His 1916 apparatus was built to measure the slope.
Both errors lie inside the $0.8\%$ standard error of
equation~\eqref{eq:se_honest}.

\subsection{The Slope Planck's Constant Predicted}
\label{sec:prediction}

Einstein derived the slope $h/e$ from Planck's radiation law.  Planck fitted
$h$ to the blackbody spectrum in 1901 and obtained $6.55\times10^{-27}$\,%
erg$\cdot$s.\footnote{Planck's 1901 value and the 1913 value in
Table~\ref{tab:comparison} both come from secondary sources and should be
checked against the originals.}  Millikan measured
$e = 1.592\times10^{-19}$\,C in his oil-drop experiment.  Those two numbers
predict a slope
\begin{equation}\label{eq:predicted_slope}
  \left(\frac{h}{e}\right)_{\text{predicted}}
  = \frac{6.55\times10^{-34}}{1.592\times10^{-19}}
  = 4.114\times10^{-15}\text{ V$\cdot$s}.
\end{equation}
The five-point sodium regression gives $4.128\times10^{-15}$\,V$\cdot$s.  The
measured slope exceeds the predicted slope by $0.35\%$, inside the $0.8\%$
standard error of \eqref{eq:se_honest}.

No photoelectric measurement entered the prediction.  Planck fitted $h$ to the
blackbody spectrum in 1901.  Einstein derived the slope in 1905.  Millikan
measured it in 1916.  The regression had no free parameter with which to meet
the prediction: Einstein fixed $\beta$ in advance and left only $\alpha$ free.
This comparison, and not the comparison with the modern $h$ in
Section~\ref{sec:comparison}, is the out-of-sample test of
equation~\eqref{eq:einstein}.

Both 1916-era constants were too small.  Planck's $h$ was low by $1.15\%$ and
Millikan's $e$ by $0.64\%$.  The errors partly cancel in the ratio.  The
predicted slope was low by $0.52\%$ and the measured slope by $0.17\%$.  The
agreement in \eqref{eq:predicted_slope} therefore rests in part on two errors
of the same sign, and a reader who wants the prediction tested against modern
constants should compare $4.128$ with $4.136\times10^{-15}$\,V$\cdot$s.

%======================================================================
\section{Maximum Likelihood Estimation of Einstein's Parameters}
\label{sec:mle}

\subsection{The Parametric Model}

Einstein's photoelectric equation defines a three-parameter
statistical model.  The free parameters are:
\begin{itemize}
  \item $\beta = h/e$ --- the universal slope (V$\cdot$s), encoding Planck's constant;
  \item $\alpha = -W/e$ --- a metal-specific intercept (V), equal to minus
    the work function in volts under the correction convention
    of~\eqref{eq:alpha_meaning};
  \item $\sigma^2$ --- the variance of the measurement noise $\varepsilon_i$.
\end{itemize}
The full parametric model for a single metal over $n$ spectral lines is:
\begin{equation}\label{eq:para_model}
  V_{0i} = \alpha + \beta\,\nu_i + \varepsilon_i, \quad
  \varepsilon_i \overset{\text{iid}}{\sim} \mathcal{N}(0,\sigma^2),
  \quad i=1,\ldots,n.
\end{equation}
The Gaussian assumption is natural: measurement noise in the electrometer
readings arises from many small independent sources (thermal Johnson noise,
vibration, residual-gas collisions), for which a central-limit argument is
compelling.

\subsection{The Log-Likelihood}

The probability density of a single observation $(V_{0i},\nu_i)$ under
model~\eqref{eq:para_model} is
\[
  p(V_{0i}\mid\nu_i;\,\alpha,\beta,\sigma) =
  \frac{1}{\sqrt{2\pi}\,\sigma}\exp\!\left(-\frac{(V_{0i}-\alpha-\beta\nu_i)^2}{2\sigma^2}\right).
\]
For $n$ independent observations the log-likelihood is
\begin{equation}\label{eq:loglik}
  \ell(\alpha,\beta,\sigma) =
  -\frac{n}{2}\ln(2\pi) - n\ln\sigma
  - \frac{1}{2\sigma^2}\underbrace{\sum_{i=1}^n(V_{0i}-\alpha-\beta\nu_i)^2}_{\displaystyle\text{RSS}(\alpha,\beta)}.
\end{equation}

\subsection{MLE Equals OLS Under Gaussian Errors}

For fixed $\sigma$, maximising~\eqref{eq:loglik} over $(\alpha,\beta)$ requires
minimising $\mathrm{RSS}(\alpha,\beta)$ --- exactly the OLS objective.
Therefore:
\begin{equation}\label{eq:mle_ols}
  \boxed{(\hat\alpha_{\mathrm{MLE}},\,\hat\beta_{\mathrm{MLE}})
         = (\hat\alpha_{\mathrm{OLS}},\,\hat\beta_{\mathrm{OLS}}).}
\end{equation}
The equivalence holds for any $\sigma$.

\subsection{MLE for the Noise Variance}

Substituting the maximising $(\hat\alpha,\hat\beta)$ into~\eqref{eq:loglik} and
optimising over $\sigma$ gives:
\begin{equation}\label{eq:sigma_mle}
  \hat\sigma_{\mathrm{MLE}}^2 = \frac{\mathrm{RSS}}{n}.
\end{equation}
The OLS estimator uses the degrees-of-freedom correction:
\begin{equation}\label{eq:sigma_ols}
  s^2 = \frac{\mathrm{RSS}}{n-2},
\end{equation}
which is unbiased ($\mathbb{E}[s^2]=\sigma^2$).  The MLE $\hat\sigma^2_{\mathrm{MLE}}$
is biased downward by the factor $(n-2)/n$; for $n=5$ the correction is $3/5$.

For the five-point sodium regression,
$\mathrm{RSS} = 1.03\times10^{-7}$\,V$^2$ (from the residuals in
Section~\ref{sec:regression}).  Hence:
\begin{align*}
  \hat\sigma_{\mathrm{MLE}} &= \sqrt{1.03\times10^{-7}/5}
  = 0.14\times10^{-3}\text{ V} = 0.14\,\text{mV}, \\
  s_{\mathrm{OLS}}  &= \sqrt{1.03\times10^{-7}/3}
  = 0.18\times10^{-3}\text{ V} = 0.18\,\text{mV}.
\end{align*}
Both lie nearly two orders of magnitude below Millikan's stated experimental
uncertainty of $\pm10$\,mV.  Sections~\ref{sec:data}
and~\ref{sec:regression} give the reason: the residuals are rounding error in
the tabulated stopping potentials.  I therefore use the instrumental value
$\sigma = 10$\,mV as the plug-in, not $s$.  The six-line fit, which includes
Millikan's one directly quoted stopping potential, gives $s = 11$\,mV.

\subsection{Fisher Information and the Cram\'er-Rao Bound}
\label{sec:fisher}

The Fisher information matrix bounds the precision attainable by any
unbiased estimator of $(\alpha,\beta)$.  For model~\eqref{eq:para_model} with
$\sigma$ known, the Fisher information is
\begin{equation}\label{eq:fisher}
  \mathcal{I}(\alpha,\beta) =
  \frac{1}{\sigma^2}\,X^\top X =
  \frac{1}{\sigma^2}
  \begin{pmatrix} n & \sum_i\nu_i \\[4pt] \sum_i\nu_i & \sum_i\nu_i^2
  \end{pmatrix},
\end{equation}
where the row $(1,\,\nu_i)$ of the $n\times2$ design matrix $X$ maps
$(\alpha,\beta)$ into the conditional mean of $V_{0i}$.
The Cram\'er-Rao bound on the variance of any unbiased estimator of $\beta$ is
\begin{equation}\label{eq:crb}
  \mathrm{Var}(\tilde\beta) \geq
  \bigl[\mathcal{I}^{-1}(\alpha,\beta)\bigr]_{\beta\beta}
  = \frac{\sigma^2}{S_{\nu\nu}},
\end{equation}
where $S_{\nu\nu}=\sum_i(\nu_i-\bar\nu)^2$.  The OLS/MLE estimator
$\hat\beta$ achieves this bound: $\mathrm{Var}(\hat\beta) = \sigma^2/S_{\nu\nu}$
(by the Gauss-Markov theorem, OLS is the Best Linear Unbiased Estimator).

For the five-point sodium regression, with $\sigma = 10$\,mV:
\begin{equation}\label{eq:crb_num}
  \mathrm{SE}_{\mathrm{CR}}(\hat\beta) =
  \frac{\sigma}{\sqrt{S_{\nu\nu}}} =
  \frac{1.0\times10^{-2}\text{ V}}{\sqrt{9.277}\times10^{14}\text{ Hz}}
  = 3.3\times10^{-17}\text{ V$\cdot$s},
\end{equation}
corresponding to a relative precision of $0.79\%$, which OLS attains exactly.
The bound depends on the design and not on the readings.  $S_{\nu\nu}$ depends
only on the spectral lines Millikan chose.  The reconstruction of
Section~\ref{sec:data} corrupts the residuals and leaves $S_{\nu\nu}$
untouched.

\paragraph{Design implications.}  By~\eqref{eq:crb},
$\mathrm{SE}(\hat\beta) \propto 1/\sqrt{S_{\nu\nu}}$.  Spreading the
frequencies further apart shrinks the standard error whatever the value of
$n$.  This justifies Millikan's choice to span nearly a factor of two in $\nu$ (5461\,\AA\ to 2535\,\AA), rather
than the narrow UV ranges, $60$--$70$\,nm wide, available to Hughes and to
Richardson \& Compton.

Table~\ref{tab:crb_comparison} quantifies the improvement.  Restrict Millikan's
sodium lines to 4047--3126\,\AA.  That $92$\,nm window is the narrowest his
line set permits and the closest analogue to the $\sim$60--70\,nm bands of
Hughes and of Richardson \& Compton.  It gives
$S_{\nu\nu} = 2.44\times(10^{14})^2$\,Hz$^2$, against $9.28$ for his
five-point design and $24.72$ for the full six.  Holding $\sigma$ fixed, the
standard error on $h/e$ falls by a factor of $\sqrt{9.28/2.44} = 1.9$ and
then $\sqrt{24.72/2.44} = 3.2$.  Widening the frequency range improved the
design.  Collecting more readings would not have.

\begin{table}[ht]
\centering
\caption{Effect of frequency range on the Cram\'er-Rao bound for
$\hat\beta = h/e$, using Millikan's sodium lines and a common instrumental
$\sigma = 10$\,mV.  ``Narrow'' uses the three lines spanning the window
closest to the ranges available to Hughes (1912) and Richardson \& Compton
(1912).  Because $\sigma$ is held fixed, the last two columns isolate the
contribution of the experimental design alone.}
\label{tab:crb_comparison}
\small
\begin{tabular}{lccccr}
\toprule
Design & $n$ & Span & $S_{\nu\nu}$ & SE($\hat\beta$) & Rel. \\
 & & (nm) & ($10^{28}$\,Hz$^2$) & ($10^{-17}$\,V$\cdot$s) & SE \\
\midrule
Narrow (4047--3126\,\AA)        & 3 & 92  & 2.44  & 6.4 & 1.55\% \\
Millikan 5-pt (5461--3126\,\AA) & 5 & 234 & 9.28  & 3.3 & 0.79\% \\
Millikan full (5461--2535\,\AA) & 6 & 293 & 24.72 & 2.0 & 0.49\% \\
\bottomrule
\end{tabular}
\end{table}

\subsection{Pooled Estimation Across Metals}
\label{sec:panel}

A more efficient estimator exploits the cross-metal restriction: $\beta = h/e$
is universal.  The pooled model is
\begin{equation}\label{eq:panel}
  V_{0im} = \alpha_m + \beta\,\nu_i + \varepsilon_{im},
  \quad m\in\{\mathrm{Na},\mathrm{Li}\},
\end{equation}
where the intercept $\alpha_m$ is specific to metal $m$ and the slope $\beta$
is common to Na and Li.  The pooled OLS slope estimator is
\begin{equation}\label{eq:pool_beta}
  \hat\beta_{\mathrm{pool}} =
  \frac{S_{\nu V}^{\mathrm{Na}} + S_{\nu V}^{\mathrm{Li}}}
       {S_{\nu\nu}^{\mathrm{Na}} + S_{\nu\nu}^{\mathrm{Li}}}.
\end{equation}
From Section~\ref{sec:regression}, $\hat\beta_{\mathrm{Na}} = 4.128$ and
$\hat\beta_{\mathrm{Li}} = 4.127$ (both $\times10^{-15}$\,V$\cdot$s).
Pooling by \eqref{eq:pool_beta} gives
\[
  \hat\beta_{\mathrm{pool}} =
  \frac{3.830 + 6.457}{9.277 + 15.647}
  = 4.127\times10^{-15}\text{ V$\cdot$s},
\]
with $S_{\nu\nu}$ raised from $9.28$ to $24.92\times(10^{14})^2$ and hence
$\mathrm{SE}(\hat\beta_{\mathrm{pool}}) = \sigma/\sqrt{24.92}\times10^{-14}
= 2.0\times10^{-17}$\,V$\cdot$s, a relative precision of $0.49\%$.  Pooling
two metals buys what adding the sixth sodium line buys.  Both enlarge
$S_{\nu\nu}$.

The Wald statistic for testing the slope-equality restriction
$H_0\colon\beta_{\mathrm{Na}} = \beta_{\mathrm{Li}}$ is
\[
  z = \frac{\hat\beta_{\mathrm{Na}} - \hat\beta_{\mathrm{Li}}}
      {\sqrt{\mathrm{SE}^2(\hat\beta_{\mathrm{Na}})
           + \mathrm{SE}^2(\hat\beta_{\mathrm{Li}})}}
  \approx \frac{0.002\times10^{-15}}
      {\sqrt{2}\times3.3\times10^{-17}}
  \approx 0.04,
\]
far below the critical value of $1.96$.  The restriction is not rejected.

The test has little power.  With a standard error of $0.8\%$ on each slope it
would fail to reject slope differences of up to about $2\%$, so it excludes
only gross violations of universality.  The test is also partly circular.
Four of the five lithium observations are the sodium observations shifted by a
constant (Section~\ref{sec:data}), which forces $\hat\beta_{\mathrm{Li}}$
toward $\hat\beta_{\mathrm{Na}}$.  The near-zero $z$ reflects that
construction.  Millikan reported parallel $V_0$--$\nu$ curves for the two
metals at his stated precision.  Nothing in the pooled data contradicts
universality.  These ten numbers do not establish it.

%======================================================================
\section{The Contact-EMF Identification Problem}
\label{sec:contact}

The contact EMF $v_c$ between the photoemitting metal and the collector is a
nuisance parameter.  In the raw readings of \eqref{eq:obs_raw} the intercept
is $-W/e-v_c$, which conflates the work function with an apparatus term that
was often of comparable size and was not stable across the early experiments.

\paragraph{Pre-Millikan experiments.}  The early experiments did not measure
or control $v_c$.  If $v_c$ drifted between measurements at different
frequencies, the apparent stopping potentials would shift inconsistently and
bias both slope and intercept.  Hughes and Richardson \& Compton
measured different lines on different days with no guarantee of constant $v_c$,
which Millikan \citeyearpar{Millikan1916} identified as a contributing cause
of their $h$ values being 12--47\% below Planck's.

\paragraph{Millikan's solution.}  Millikan measured all six spectral lines
in a single continuous session after a freshly shaved surface had stabilised
(typically 1--2 hours after shaving).  He verified that $v_c$ remained constant
to $\pm 0.01$\,V over the session by using the Kelvin double-balance.  When
$v_c$ is constant across all observations, it enters as a pure intercept shift
and does not affect $\hat\beta$.  The slope $\hat\beta = h/e$ is then estimated
without bias regardless of the unknown value of $v_c$.

\paragraph{Separate identification of $W$.}  The independent Kelvin
measurement $v_c = 2.51$\,V is what makes the intercept interpretable at all.
Having used it to convert his raw readings into true stopping potentials,
Millikan could read the work function directly off the fitted intercept
$\hat\alpha = -1.813$\,V of equation~\eqref{eq:alphahat_num}, since by
\eqref{eq:alpha_meaning}
\begin{align*}
  W_{\mathrm{Na}}/e &= -\hat\alpha = 1.81\text{ eV},\\
  W_{\mathrm{Na}} &= 1.81\times1.602\times10^{-19}
    = 2.90\times10^{-19}\text{ J}.
\end{align*}
The corresponding lithium figure is $-\hat\alpha_{\mathrm{Li}} = 2.35$\,eV.
The dependent variable already contains $v_c$.  Adding it again double-counts
it.

These estimates are the weakest quantities in the analysis.  The modern value
for clean bulk sodium is $2.36$\,eV, and the intercept-implied figure is low
by about $0.5$\,eV.  Estimation error accounts for part of the gap.  The
intercept extrapolates to $\nu=0$, far outside the observed range
($5.5$--$11.8\times10^{14}$\,Hz), and its standard error is
$\sigma\sqrt{1/n + \bar\nu^2/S_{\nu\nu}} \approx 0.08$\,V at
$\sigma = 10$\,mV.  Surface physics accounts for more.  Millikan's shaved
surfaces sat in $10^{-2}$\,mmHg of residual gas, and adsorbed gas shifts a
metal's work function by several tenths of a volt.  $W$ varies by some
$0.3$\,eV across the crystal faces of a single metal.  Sodium has no one work
function for the intercept to be compared against.

The two coefficients are identified on different terms.  The slope is a
universal constant.  Variation within a run identifies it, and any disturbance
common to all lines leaves it alone.  The intercept is a surface-specific
quantity.  Extrapolation identifies it, and the conditions Millikan could not
control contaminate it.  Einstein's equation is better identified in its slope
than in its intercept.  The slope carries the physics.

%======================================================================
\section{Nine Two-Point Determinations}
\label{sec:twopoint}

The five- and six-point regressions cannot supply a credible standard error
from their own residuals (Section~\ref{sec:data}).  Millikan's Table~II can.
He estimated the slope from nine pairs of spectral lines, always pairing one
line with 5461\,\AA, chosen for its high intensity and its precise balance
point.  Each row is a separate session with a freshly shaved surface.  The
nine values are replications.  Their dispersion measures what varied from run
to run: surface condition, residual gas, drift in $v_c$, reading error.

\begin{table}[ht]
\centering
\caption{Millikan (1916), Table~II: nine two-point slope determinations
for sodium.  Each row is an independent measurement session with a fresh
surface shaving.  The mean and standard deviation are computed from these
nine observations.}
\label{tab:two_point}
\begin{tabular}{ccc}
\toprule
Session & $\lambda$ paired with 5461\,\AA & $\hat\beta$ ($10^{-15}$\,V$\cdot$s) \\
\midrule
1 & 3126\,\AA & 4.11 \\
2 & 3650\,\AA & 4.14 \\
3 & 3126\,\AA & 4.10 \\
4 & 3650\,\AA & 4.12 \\
5 & 3126\,\AA & 4.24 \\
6 & 4047\,\AA & 3.98 \\
7 & 2535\,\AA & 4.04 \\
8 & 3126\,\AA & 4.24 \\
9 & 4047\,\AA & 4.21 \\
\midrule
Mean   & & $\mathbf{4.131}$ \\
Std.~dev.\ ($n=9$) & & $0.087$ \\
Std.~error of mean & & $0.029$ \\
95\% CI & & $[4.06,\; 4.20]$ \\
\bottomrule
\end{tabular}
\end{table}

\noindent The confidence interval is Student-$t$ with eight degrees of
freedom ($t_{.975,8} = 2.306$); the normal approximation would give
$[4.073,\,4.189]$.

The two-point mean $4.131\times10^{-15}$\,V$\cdot$s agrees with the five-point
OLS estimate $4.128\times10^{-15}$\,V$\cdot$s to better than $0.1\%$.  The
reconstruction of Section~\ref{sec:data} has not distorted the slope.  The
session-to-session standard deviation of $0.087$, or $2.1\%$ of the mean, is
the scale of the experiment's variability.  It exceeds the $0.18$\,mV residual
scatter of the five-point fit by a factor of about a hundred.  The standard
error of the mean, $0.029\times10^{-15}$ or $0.71\%$, estimates the precision
of Millikan's determination from repeated measurement alone.

That figure sits on the Cram\'er-Rao bound of Section~\ref{sec:fisher}, which
the design and $\sigma = 10$\,mV give as $0.79\%$ for the five-line span.
Replication and the information matrix agree to within a tenth of a percent.
This agreement warrants the uncertainty attached to $h$ in this paper.  The
$1$--$4\%$ spread of individual sessions is what $\sigma\approx10$\,mV applied
to a two-point slope produces.

%======================================================================
\section{Comparison of Historical Estimates}
\label{sec:comparison}

Table~\ref{tab:comparison} compares the estimates of $h$ from all major
pre-1920 photoelectric experiments.

\begin{table}[ht]
\centering
\caption{Historical estimates of Planck's constant $h$ from photoelectric
experiments.  Millikan's attribution of earlier errors is discussed in
\citet[pp.\ 373--375]{Millikan1916}.  For the Millikan rows the slope
$\hat\beta$ is converted to $h$ twice: once with the charge Millikan himself
had measured, $e = 1.592\times10^{-19}$\,C, and once with the modern
$e = 1.602177\times10^{-19}$\,C.  Errors are relative to the defined value
$h = 6.626070\times10^{-27}$\,erg$\cdot$s.  The two Planck values come from
secondary sources and should be checked against the originals.}
\label{tab:comparison}
\small
\begin{tabular}{p{3.6cm}rcrrr}
\toprule
Authors & Year & Metal(s) & $h$, Millikan's $e$ & Err.
  & $h$, mod. $e$ \\
\midrule
Hughes        & 1912--13 & various & 3.55--5.85 & $-$12 to $-$46\% & --- \\
Richardson \& Compton & 1912 & 8 metals & 3.5--5.8 & $-$12 to $-$47\% & --- \\
Millikan (Na, 5-pt)  & 1916 & Na & 6.573 & $-$0.81\% & 6.615 \\
Millikan (Na, 9 two-pt)& 1916 & Na & 6.577 & $-$0.74\% & 6.619 \\
Millikan (Li, 5-pt)  & 1916 & Li & 6.570 & $-$0.85\% & 6.612 \\
Millikan (Na+Li pooled)& 1916 & Na+Li & \textbf{6.571} & $\mathbf{-0.83\%}$
  & \textbf{6.613} \\
Planck (black body)  & 1901 & --- & 6.55 & $-$1.15\% & --- \\
Planck (black body)  & 1913 & --- & 6.415 & $-$3.19\% & --- \\
\midrule
Modern (defined)     & 2019 & --- & 6.626070 & 0 & 6.626070 \\
\bottomrule
\end{tabular}
\end{table}

Millikan's photoelectric determinations sit $0.7$--$0.9\%$ below the modern
value.  His oil-drop charge, too small by $0.64\%$, accounts for almost all of
that gap.  Rerun with the modern $e$, every one of his slopes lands within
$0.22\%$ of the defined $h$, inside the $0.8\%$ standard error of the design.
His photoelectric measurement was more accurate than the number he published.

The pre-Millikan estimates are biased downward by $12$--$47\%$.  The
discrepancy exceeds any plausible sampling error by two orders of magnitude,
so it is systematic.  Three mechanisms produce it.  Stray short-wavelength
light raises the apparent $V_0$ at long-wavelength lines and flattens the
slope.  A restricted frequency range inflates the standard error, as
Table~\ref{tab:crb_comparison} quantifies, and leaves the estimate at the
mercy of a single bad line.  Contact EMF that drifts between sessions
contaminates observations individually instead of shifting a common intercept,
and biases the slope in an unpredictable direction.  Millikan's engineering
addressed all three: fresh surfaces, filtered lines, a wide frequency range,
and continuous Kelvin monitoring of $v_c$.  Only the second is a statistical
problem.  Estimation theory cannot repair the other two after the fact.

%======================================================================
\section{Millikan's Reluctance}
\label{sec:millikan_reluctance}

Millikan's 1916 paper confirmed Einstein's equation~\eqref{eq:einstein} to
to better than $1\%$.  Yet Millikan wrote:
\begin{quote}
  ``Despite then the apparently complete success of the Einstein equation,
  the physical theory of which it was designed to be the symbolic expression
  is found so untenable that Einstein himself, I believe, no longer holds
  to it.'' \citep[p.\ 388]{Millikan1916}
\end{quote}
The light-quantum hypothesis was, in Millikan's judgment, physically
implausible because it conflicted with the wave theory of diffraction and
interference.  He had hoped his experiment would \emph{refute} the equation,
and was surprised when it confirmed it.

Compton scattering \citep{Compton1923} supplied, in 1923, the independent
evidence for the particulate character of light that the photoelectric
regression could not.  Millikan came to describe his own result in these
terms:
\begin{quote}
  ``This work resulted, contrary to my own expectation, in the first direct
  experimental proof \ldots of the complete validity of the Einstein
  equation.''
\end{quote}
\noindent The sentence belongs to Millikan's later retrospective writing, not
to the Nobel lecture he delivered in May 1924.  \citet{Holton1999} and
\citet{Pais1982} trace the trajectory of his views.

A macroeconomist likewise can confirm a reduced-form predictive relationship
without endorsing the structural model that motivated it.

%======================================================================
\section{Einstein as Machine Learner}
\label{sec:ml}

\subsection{Overview}

The vocabulary of modern machine learning---supervised learning, hypothesis
class, loss function, inductive bias, feature engineering, transfer
learning---did not exist in 1905.  The operations Einstein performed are the
ones that statistical learning theory formalises.
\citet{sargent2025sources} argues that this is the rule rather than the
exception, and that the techniques now gathered under artificial intelligence
have long histories inside the sciences that used them.  This section states
the correspondence for the photoelectric case and its limits.

\subsection{Supervised Learning and Labeled Data}

The photoelectric dataset is a standard supervised regression problem.
The \emph{input feature} is the frequency~$\nu$ of the incident light;
the \emph{label} is the stopping potential~$V_0$.  The task is to learn
a function $f:\nu\mapsto V_0$ that generalises to unobserved frequencies.

Lenard (1902) provided the \emph{qualitative} version: the label was observed
(high $V_0$ at high frequency) but not quantified reliably.  Millikan (1916)
provided the \emph{quantitative} training set: six labelled pairs
$(\nu_i, V_{0i})$ for sodium and five for lithium, with precision sufficient
to estimate the slope $h/e$ to under $1\%$.

\subsection{The Hypothesis Class and Inductive Bias}
\label{sec:hyp_class}

Einstein's hypothesis class is:
\begin{equation}\label{eq:hyp_class}
  \mathcal{H} = \bigl\{f_{\alpha,\beta}:\nu\mapsto\alpha+\beta\nu
  \bigm|\alpha\in\mathbb{R},\;\beta>0\bigr\},
\end{equation}
together with the \emph{cross-metal constraint} that $\beta$ is the same
for all metals.  The class is extremely tight.  For a single metal it has
two free parameters, $(\alpha,\beta)$.  Across $M$ metals the constraint bites:
instead of $2M$ free parameters there are $M+1$---one intercept apiece and a
single shared slope---so the class is a family of parallel lines whose common
gradient is the one physical unknown.

The restriction to $\mathcal{H}$ was not ad hoc.  It followed deductively
from the quantum hypothesis $E=h\nu$ and energy conservation for photoemission.
The quantum hypothesis \emph{is} the inductive bias: it restricts the
hypothesis class from all possible functions of $(\nu, I, \lambda, \ldots)$
to a specific affine form in $\nu$ alone.  Without this bias five data points
are hopelessly insufficient: infinitely many smooth functions pass through
five points.

\subsection{Feature Engineering: Choosing $\nu$ Over Intensity
and Wavelength}

Learning begins with a choice of input features.
Lenard's experiments performed a feature importance analysis.  He
measured $V_0$ as a function of both frequency (colour) and intensity~$I$.  His
result was unambiguous: $V_0$ varies strongly with $\nu$ but not at all with~$I$.

Classical wave theory predicted the opposite: wave amplitude $A \propto
\sqrt{I}$ should determine energy delivered to an electron, so $V_0$ should
grow with $I$.  The empirical irrelevance of $I$ was the anomaly that
motivated Einstein's model.  Lenard's experiment eliminated $I$ from the
feature set and selected $\nu$ before any parametric fitting was done.

\subsection{The Loss Function and Training}

Given the hypothesis class~\eqref{eq:hyp_class} and the feature $\nu$, choose
$(\alpha,\beta)$ to minimise the empirical risk under squared-error loss:
\[
  \hat{R}(\alpha,\beta) = \frac{1}{n}\sum_{i=1}^n(V_{0i} - \alpha - \beta\nu_i)^2.
\]
This is identical to the MSE loss used in neural-network regression.  The
minimiser of $\hat{R}$ is the OLS/MLE estimator derived in
Sections~\ref{sec:regression}--\ref{sec:mle}: training is the closed-form
solution $\hat\beta = S_{\nu V}/S_{\nu\nu}$.  For Einstein's linear model,
the loss surface is a convex paraboloid in $(\alpha,\beta)$ with a unique
global minimum; no stochastic gradient descent is needed.

\subsection{Generalisation and Transfer Learning}

\paragraph{Generalisation.}
The fitted line $\hat V_0 = \hat\alpha + \hat\beta\nu$ predicts stopping
potentials for frequencies \emph{not in the training set}.  The residuals of
the five-point fit cannot establish this, since they measure rounding
(Section~\ref{sec:data}).  Two other results can.  The six-point fit, which
includes a line withheld from the five-point design, leaves a residual
standard deviation of $11$\,mV at that line.  The nine two-point sessions of
Section~\ref{sec:twopoint} each fit two lines and agree with the five-point
slope to $0.1\%$.  In
learning-theoretic terms, the relevant capacity measure for a real-valued
class is the pseudo-dimension (the Vapnik-Chervonenkis dimension is defined
for classifiers); for affine functions of a single input it equals~2, the
number of free parameters.  Five training points against a capacity of two
leave three residual degrees of freedom---enough both to fit and to assess the
fit, which is precisely what a capacity of five or more would forbid.

\paragraph{Transfer learning.}
A model trained on sodium data successfully predicts lithium stopping
potentials up to an intercept shift.  The slope $\hat\beta = h/e$ is
transferred exactly from one metal to another.  Training on sodium provides a
pre-trained feature extractor, the slope; the lithium task fine-tunes the
intercept $\alpha_{\mathrm{Li}}$, the metal-specific work-function parameter.  The Wald test of
Section~\ref{sec:panel} confirms that the zero-shot transfer of the slope
from Na to Li is not rejected by the data ($z = 0.04$), subject to the two
cautions recorded there.

\subsection{Model Selection: Quantum vs.\ Classical Wave Theory}

At the time of Einstein's paper two models were in competition:
\begin{itemize}
  \item \textbf{Classical energy-transfer model.}  The energy delivered to an
    electron is set by the wave amplitude $A \propto \sqrt{I}$ accumulated over
    the illumination time.  $V_0$ should therefore rise with intensity, rise
    with exposure, and depend on frequency only weakly and through no
    prescribed functional form.
  \item \textbf{Einstein's model.}  $V_0 = (h/e)\nu - W/e$: exactly linear in
    $\nu$, wholly independent of $I$, with a slope fixed \emph{a priori} at
    $h/e$ and common to every metal.
\end{itemize}
On these data the second model wins by any criterion.  The data bear on three
claims, and the evidence for them differs.

Lenard settled the irrelevance of intensity before anyone ran a regression,
and that alone refutes the classical energy-transfer picture.  The regression
establishes linearity in $\nu$.  The numerical value of the slope is the
sharpest test.  Planck's blackbody constant fixed $h/e$ in advance, Einstein's
prediction had no free parameter to adjust, and the measured slope came within
$0.35\%$ of it (Section~\ref{sec:prediction}).  A model-selection criterion
that scores in-sample fit credits a correct prediction and a lucky post-hoc
fit alike, and misses this last point.

The data do not select the light-quantum mechanism.
Equation~\eqref{eq:einstein} follows from energy conservation together with
quantised \emph{atomic} energy levels, and a semiclassical model in which the
radiation field remains a classical wave delivers it too
(Appendix~\ref{app:lamb_scully}).  The regression identifies the best-fitting
member of a hypothesis class and rules out a rival class.  It cannot
discriminate between two physical stories that imply the same class.
Photon-number statistics discriminated between them sixty years later.  Model
selection has this limitation generally, and Millikan's execution of it is not
at fault.

\subsection{What a Naive Machine Learning System Would Miss}

A purely data-driven system, given only the table of $(\nu_i, V_{0i})$
pairs, would learn the following.

\begin{enumerate}
  \item A linear regression package would correctly recover
    $\hat\beta = 4.128\times10^{-15}$\,V$\cdot$s.
  \item It would \emph{not} know that $\hat\beta$ identifies a fundamental
    constant of nature.  The number $4.128\times10^{-15}$ is physically
    meaningless without the quantum hypothesis that gives it the
    interpretation $h/e$.
  \item It would \emph{not} impose the cross-metal universality
    restriction $\beta_{\mathrm{Na}} = \beta_{\mathrm{Li}}$ from first
    principles.  The universality is a prediction of the quantum hypothesis;
    from a purely data-driven perspective, it is just a coincidence to be
    discovered, not a constraint to be imposed.
  \item It would have no way to decompose the intercept into the physical
    work function $-W/e$ and the apparatus nuisance $v_c$, since the two are
    not separately identified without the external Kelvin measurement
    (Section~\ref{sec:contact}).  Identification failures of this kind are
    invisible to a fitting procedure: the regression runs, reports an
    intercept, and says nothing about what the intercept means.
\end{enumerate}
Section~\ref{sec:dnn} pursues the comparison to a system with far greater
capacity than a linear regression package, and reaches the same conclusion by
a different route.

Einstein supplied the theoretical framework that gives the regression
coefficients their physical meaning.  Specifying the hypothesis class required
more intellectual effort than fitting the line.

\subsection{The Deductive--Inductive Loop}

Modern discussions of ``AI for science'' often describe a purely inductive
paradigm: given enough data, learn the law.  The photoelectric episode follows
the hypothetico-deductive loop that has driven physical science since Newton.

\begin{table}[ht]
\centering
\caption{The hypothetico-deductive loop in the photoelectric episode,
translated into the language of modern machine learning.}
\label{tab:loop}
\begin{tabular}{p{4cm}p{4cm}p{4cm}}
\toprule
Step & Scientific name & ML name \\
\midrule
Propose $E = h\nu$ & Theoretical hypothesis & Model architecture \\
Derive $V_0 = (h/e)\nu - W/e$ & Deduction &
  Hypothesis class \\
Choose $\nu$ not $I$ & Feature selection & Feature engineering \\
Measure $(\nu_i,V_{0i})$ & Experiment & Training data \\
Min.\ $\sum(V_{0i}-\hat V_{0i})^2$ & Least squares &
  MSE training \\
$|\hat\beta - h/e| < 1\%$ & Confirmation &
  Test accuracy \\
Na and Li: same $\hat\beta$ & Universality & Zero-shot transfer \\
\bottomrule
\end{tabular}
\end{table}

Five training points estimate $h$ to 0.1\% because the physics prescribes the
functional form, the number of free parameters $(\alpha,\beta)$, and the units
of the slope.  A purely inductive system that must discover the functional form
from data faces a far larger sample-complexity problem.

%======================================================================
\section{A Deep Neural Network on the Photoelectric Data}
\label{sec:dnn}

\subsection{Two Hypothesis Classes Compared}

Einstein's model and a deep neural network (DNN) represent opposite
extremes of the bias-variance spectrum for the photoelectric regression.
Einstein's hypothesis class is equation~\eqref{eq:hyp_class}: two free
parameters per metal, one shared slope across metals, pseudo-dimension~2.
The Gauss--Markov theorem and the Cram\'er--Rao bound together guarantee
that OLS achieves minimum variance among all unbiased estimators;
five training points are more than sufficient.

A DNN with $L$ hidden layers and $N$ neurons per layer is a universal
function approximator: for any continuous target function on a compact
interval, there exists a network that approximates it to arbitrary
precision \citep{Hornik1989}.  For a one-dimensional input ($\nu$) and
one-dimensional output ($V_0$), even a modest architecture---three
hidden layers with 16 ReLU neurons each---has approximately
\[
  p = (1\cdot16+16)+(16\cdot16+16)+(16\cdot16+16)+(16\cdot1+1)
    = 593
\]
free parameters.  For ReLU networks the pseudo-dimension grows like
$p\log_2 p$, here $593\log_2 593 \approx 5{,}500$ \citep{Vapnik1998}.

The literature offers no single conversion from a capacity figure to a
required sample size.  The uniform-convergence bound $n = O(d/\epsilon^2)$,
with $d$ the pseudo-dimension, gives $n \sim 10^7$ for $\epsilon = 1\%$.  The
bound is vacuous at that size.  The practitioner's rule $n \gtrsim 10p$ gives
$n \gtrsim 6{,}000$.  I use the rule throughout as the ``sufficient $n$'' for
the DNN.  The rule is a heuristic.  The formal bound is four orders of
magnitude worse.

Capacity bounds of this kind are loose.  Overparameterised networks generalise
in regimes where uniform-convergence theory predicts they cannot
\citep{Belkin2019}, and the argument below does not rest on a claim that a DNN
provably fails.  The difficulty at $n=5$ concerns identifiability and not
generalisation error.  Section~\ref{sec:epochIII} makes it precise.  The
comparison with Einstein's two-parameter model turns on degrees of freedom,
not on a VC bound.

\subsection{Epoch I---Lenard's Qualitative Findings (1902)}

Lenard established two qualitative facts: stopping potential rises with
frequency and is independent of intensity.  Lenard recorded no
$(\nu_i,\,V_{0i})$ pairs.  A DNN requires
labelled numerical examples;
with none, it cannot begin training.  Einstein required zero labelled
examples: the quantum hypothesis plus energy conservation yielded
equation~\eqref{eq:einstein} as a deduction.  Einstein's ``model fitting'' was
a one-paragraph derivation.  The network's fitting cannot begin.

\subsection{Epoch II---Hughes and Richardson \& Compton (1907--1914)}

The intermediate experiments produced roughly 25--40 numerical
$(\nu_i,\,V_{0i})$ pairs, clustered in a narrow 60\,nm UV range with
uncontrolled contact EMFs and stray scattered light.
Three problems foreclose reliable DNN learning.

\paragraph{Underdetermination.}
With $n \approx 30$ and $p = 593$, the DNN is massively overparameterised.
Gradient descent finds a minimum-norm interpolant that passes through
all 30 training points exactly (zero training loss), but is effectively
unconstrained between them.  Applying $\ell_2$ weight decay suppresses
oscillations, but the optimal regularisation strength cannot be selected
without a held-out validation set; with $n=30$ this is statistically
isolated.

\paragraph{Restricted frequency range.}
A linear, a quadratic, and a square-root form fit almost equally well over a
$60$--$70$\,nm UV band ($\Delta\nu \approx 1.5\times10^{14}$\,Hz).  The DNN has
no mechanism
to infer which of them is the global functional form.

\paragraph{Systematic label contamination.}
Uncontrolled contact EMFs imply that the effective label for observation
$i$ is $V_{0i} + \delta_i$, where $\delta_i$ is unknown and
observation-specific.  The DNN receives contaminated targets and,
lacking any physical model of the contaminant, cannot separate the
linear physical trend from the systematic noise.  The learned function
absorbs both, producing a biased and uninterpretable output.

\subsection{Epoch III---Millikan's Data (1916)}
\label{sec:epochIII}

Millikan's sodium experiment was the cleanest of the era.  It yielded five
usable $(\nu_i, V_{0i})$ pairs for the primary determination.

\paragraph{The interpolation regime.}
With $n=5$ and any DNN with $p>5$ (which includes every
useful architecture), the network operates in the interpolation
regime: it memorises all five training points exactly,
producing zero training loss \citep{Belkin2019}.  The resulting
function is unconstrained everywhere between the training points.
Even the best-case minimum-norm interpolant---``well-behaved'' in the
double-descent sense---provides no basis for slope inference: there
are no residual degrees of freedom with which to estimate the noise
variance or a standard error on the slope.

\paragraph{Loss of inferential machinery.}
OLS with $n=5$ and $p=2$ free parameters uses $n-2=3$ residual
degrees of freedom to estimate $\hat\sigma^2$ and hence
$\mathrm{SE}(\hat\beta)$.  Any DNN with more than three parameters
exhausts those degrees of freedom entirely.  The DNN cannot report
a confidence interval for $h/e$, cannot test whether the slope
differs from zero, and cannot quantify the 0.1\% relative precision
that OLS achieves at the Cram\'er--Rao bound.

\paragraph{No basis for universality.}
The universality of the slope $\beta = h/e$ across metals is a
theoretical prediction of the quantum hypothesis, not an empirical
feature discoverable from sodium data alone.  A DNN trained on
sodium would need to be retrained from scratch on lithium.  The
only mechanism by which the slope could transfer automatically from
Na to Li is an explicit architectural constraint that encodes the
universality---which is equivalent to importing the physics into
the model.

\subsection{The Modern Era: Large Datasets}

By the mid-twentieth century, photoelectric measurements spanned
decades of metals, surface preparations, and photon energies.
Pooling even a modest subset---say $n=5{,}000$ clean observations
over two octaves of frequency for a single metal---brings a DNN into the
classical regime ($n\gg p$), if only barely on the reckoning above.

With smooth activations and appropriate regularisation, the
DNN would converge, as $n\to\infty$, to the
minimum-population-risk function in its hypothesis class.  Because
the true function is exactly linear, the minimum-risk function in
any sufficiently rich class \emph{is} the linear function.  The
DNN would eventually recover the line.  It would still not identify the
constant, impose the cross-metal restriction, decompose the intercept, or
match Einstein's sample size:
\begin{itemize}
  \item Identify the learned slope $\hat{w} \approx
    4.128\times10^{-15}$\,V$\cdot$s as the ratio $h/e$ of two
    independently measured constants.  The slope appears as a
    numerical weight, not as a ratio of fundamental quantities.
  \item Impose the cross-metal constraint $\beta_{\mathrm{Na}} =
    \beta_{\mathrm{Li}} = \cdots$ from first principles.  Universality
    would appear as an empirical pattern to be confirmed, not a
    structural prior to be imposed.
  \item Decompose the intercept into a physical work function $W/e$
    and an apparatus nuisance $v_c$ without independent
    measurements of $v_c$.
  \item Reduce the required sample from $O(10^3)$ examples to five.
    A correctly specified two-parameter linear model needs five points;
    no architectural innovation can replicate this compression without
    importing the physics.
\end{itemize}

Modern symbolic regression systems such as AI Feynman \citep{Udrescu2020}
can, in principle, rediscover the linear form and even match $\hat{w}$ to
$h/e$ if given units and physical constants as hardcoded inputs.  But
that hardcoding supplies prior physical knowledge.  The $(\nu,V_0)$ data do
not contain it.

\subsection{Summary: Sample Complexity and the Value of Theory}

Table~\ref{tab:sample_complexity} quantifies the contrast.

\begin{table}[ht]
\centering
\caption{Sample-complexity comparison between Einstein's two-parameter
linear model (pseudo-dimension $=2$, derived from the quantum hypothesis) and
a three-layer DNN with 16 neurons per layer ($p=593$ parameters,
pseudo-dimension $\approx 5{,}500$).  ``Sufficient $n$'' uses the
practitioner's rule $n\gtrsim10p \approx 6{,}000$; the formal
uniform-convergence requirement $O(d/\epsilon^2)$ is some four orders of
magnitude larger and is not used here.}
\label{tab:sample_complexity}
\begin{tabular}{p{2.8cm}rp{2.4cm}p{2.4cm}p{4.0cm}}
\toprule
Epoch & $n$ available & Einstein model & DNN (VC$\approx$5700) & DNN outcome \\
\midrule
Lenard 1902   & $\approx 0$  & Sufficient$^\dagger$ & Insufficient
  & Cannot train (no labels) \\
Hughes 1912   & $\sim 30$    & Sufficient
  & Insufficient ($\times$200) & Memorises; no generalisation \\
Millikan 1916 & 5            & Sufficient
  & Insufficient ($\times$1200) & Interpolates; no inference \\
Modern era    & $\sim 5000$  & Sufficient
  & Short by $\times$1.2 & Recovers linear shape;
  cannot identify $h/e$ \\
\bottomrule
\multicolumn{5}{p{14cm}}{$^\dagger$ The functional form is derived from
  theory; only two free parameters remain to be estimated from data.}
\end{tabular}
\end{table}

DNNs are good learners for high-dimensional problems in which no
parsimonious parametric family is known: image recognition, protein
structure prediction, language modelling.  The photoelectric problem sits at
the opposite extreme: a one-dimensional input, a one-dimensional output, and
five to thirty training examples.

Einstein's proposal $E=h\nu$ reduced the hypothesis space from the
infinite-dimensional function class of a DNN to a two-dimensional manifold of
parallel lines, and with it the required number of training examples by three
orders of magnitude.  He did so before a single labelled example was
available.

%======================================================================
\section{Conclusions}
\label{sec:conclusions}

Einstein's photoelectric equation is a single-regressor linear model with
one slope ($h/e$) and one intercept ($-W/e$, once the contact EMF has been
measured and removed).  The regression is trivial
given clean data.  Obtaining clean data was the central challenge.  Fresh
metal surfaces free of oxide eliminate work-function noise.  High vacuum
eliminates atmospheric contamination.  Spectrally pure light eliminates bias
from stray UV.  A wide frequency range reduces the standard error of the
slope.  Control of $v_c$ identifies the intercept.

Millikan's decade-long experimental program was a campaign to
raise data quality until the OLS standard errors were small enough to
distinguish a specific value of the slope ($h/e = 4.136\times10^{-15}$\,V$\cdot$s)
from the alternatives implied by earlier, discrepant theories.

The final five-point sodium regression gives
\begin{align*}
  \hat\beta &= 4.128\times10^{-15}\text{ V$\cdot$s} \;\pm\; 0.8\%
    \quad (h = 6.62\times10^{-34}\text{ J$\cdot$s at the modern }e,
    \text{ }0.17\%\text{ low}), \\
  \hat\alpha &= -1.813\text{ V}
             \quad (W_{\mathrm{Na}}/e = 1.81\text{ eV}).
\end{align*}
The quoted standard error is the Cram\'er-Rao bound at Millikan's
instrumental precision of $10$\,mV, not the residual-based figure, which the
reconstruction of the tabulated potentials renders meaningless
(Section~\ref{sec:data}).  It agrees with the $0.7\%$ standard error of his
nine independent two-point sessions, and with his own claim of ``not more
than $0.5\%$'' for the full six-line design.

Five observations sufficed to measure a fundamental constant to under one
percent.  The physics had already fixed the functional form, the number of
free parameters, and the units of the slope.  A learner obliged to discover
those from data would need thousands of observations and would still not know
what it had found.  Theory compresses the sample size that data must supply.

Millikan's fit selected Einstein's functional form over its classical rival.
It could not select the light-quantum mechanism.  The same equation follows
from a semiclassical model in which the field is never quantised
(Appendix~\ref{app:lamb_scully}).  Photon-antibunching experiments settled
that question sixty-one years later (Appendix~\ref{app:antibunching}).  No one
would have performed them without the quantum theory of optical coherence to
say which observable mattered.  Theory sets the hypothesis class.  Theory then
tells the experimenter which measurement discriminates within it.

%======================================================================
\appendix

\section{Einstein's Equation in Modern Physics}
\label{app:modern}

\subsection{The Equation Remains Exact}

Einstein's photoelectric equation~\eqref{eq:einstein} has not been
superseded.  The equation states energy conservation for the fastest
photoelectron, the electron that starts at the Fermi level and reaches the
detector without inelastic loss.  It survives intact in the
quantum-electrodynamical treatment, where single-photon photoemission is a
first-order process whose energy balance is exactly \eqref{eq:einstein}.

The word applies to the energy balance and not to the measured spectrum.  A
real metal broadens the photoemission edge over $\sim k_BT$, so the stopping
potential locates an edge and not a sharp cutoff.  $W$ is a property of the
crystal face and of the adsorbate layer, not of the element.  Final-state
relaxation shifts core-level energies by amounts that matter at spectroscopic
precision.  The
work function $W$ is the energy required to move an electron from the Fermi
level to the vacuum level, and density-functional theory computes it to a few
tenths of an eV for clean surfaces.  Millikan's regression tests the energy
balance, which is exact; the residual scatter he saw came from the surface
physics, which is not.

The equation is exploited continuously in surface science.  Ultraviolet
Photoelectron Spectroscopy (UPS) and Kelvin probe instruments apply
equation~\eqref{eq:einstein} to measure work functions
of semiconductor films, organic photovoltaic layers, and heterogeneous
catalysts, routinely resolving differences of a few meV.

\subsection{X-ray Photoelectron Spectroscopy}
\label{app:xps}

When the incident photon energy $h\nu$ is drawn from the X-ray range
(100\,eV--10\,keV), the emitted electron's kinetic energy $eV_0$
identifies the \emph{binding energy} $E_B$ of the core electron level
from which it was ejected \citep{Siegbahn1967}:
\begin{equation}\label{eq:xps}
  E_B = h\nu - eV_0 - \phi_{\mathrm{spec}},
\end{equation}
where $\phi_{\mathrm{spec}}$ is the spectrometer work function, a
calibration constant.  Equation~\eqref{eq:xps} is
equation~\eqref{eq:einstein} rearranged: the unknown is now the binding
energy rather than $h/e$.  Because each element has a unique set of
core-level binding energies, X-ray Photoelectron Spectroscopy (XPS/ESCA)
provides quantitative, element-specific chemical analysis of surfaces.
Kai Siegbahn received the 1981 Nobel Prize in Physics for this technique,
which is now one of the most widely used methods in materials science,
semiconductor device fabrication, and surface catalysis.

In statistical terms, XPS is a multivariate extension of Millikan's
regression: dozens of binding energies are estimated simultaneously by
inverting Einstein's linear equation at multiple photon energies.

\subsection{Multi-Photon Photoelectric Effect}
\label{app:multiphoton}

Einstein's equation assumes each emitted electron absorbs exactly one
photon.  At the peak intensities of pulsed lasers
($I \gtrsim 10^{10}$\,W/cm$^2$), an electron can sequentially absorb
$n > 1$ photons before escaping the surface.  The $n$-photon
generalisation of equation~\eqref{eq:einstein} is
\begin{equation}\label{eq:multiphoton}
  eV_0 = n\,h\nu - W, \qquad n = 1, 2, 3, \ldots
\end{equation}
Above-threshold ionisation (ATI)---photoelectron emission at kinetic
energies exceeding the single-photon threshold by integer multiples of
$h\nu$---was first observed by \citet{Agostini1979}.  The resulting
photoelectron spectrum displays a comb of peaks separated by $h\nu$,
each satisfying equation~\eqref{eq:multiphoton} for the corresponding
$n$.  ATI extends hypothesis class~\eqref{eq:hyp_class} to allow non-unit photon
multiplicity.  The $n$-photon branch has slope $nh/e$, not $h/e$: the class
becomes a family of lines whose gradients are integer multiples of a single
universal constant, with the metal still entering only through the common
intercept.  In regression terms this is a mixture model over a discrete latent
index $n$.  The gradients are quantised rather than free, and a single
parameter governs every branch.

\subsection{Does the Photoelectric Effect Prove the Quantisation of Light?}
\label{app:lamb_scully}

\citet{Lamb1969} showed that the photoelectric stopping-potential equation
can be derived from a \emph{semiclassical model}: the electromagnetic
field is treated as a classical wave while only the atom (electron) is
quantised.  In this model a classical wave drives transitions between
quantised atomic energy levels, and the ejected-electron kinetic energy
still satisfies equation~\eqref{eq:einstein}.  The reason is that
equation~\eqref{eq:einstein} follows from energy conservation and the
quantised \emph{atomic} level structure alone; it does not require the
radiation field itself to be quantised.

Millikan's regression confirms the functional form $V_0 \propto \nu$ and
measures $h/e$ to better than $1\%$.  It does not distinguish the quantised-field
hypothesis from semiclassical alternatives.  The slope $h/e$ gives Planck's
constant whether or not the field is composed of quanta.

\subsection{Photon Antibunching: Definitive Evidence for Field Quantisation}
\label{app:antibunching}

The genuine experimental proof that the radiation field must be quantised
came not from the photoelectric effect but from photon-counting
experiments in quantum optics.  The relevant observable is the
statistical regularity---or irregularity---with which successive photons
are emitted.

\citet{HanburyBrown1956} observed \emph{photon bunching} in thermal
(chaotic) light: photon pairs from a thermal source arrive at two
detectors with a coincidence rate exceeding a Poisson benchmark, reflecting
the Bose-Einstein statistics of thermal photons.  \citet{Glauber1963}
developed the quantum theory of optical coherence, introducing coherent
states of the radiation field (the states most closely resembling classical
waves) and deriving the photon-number statistics that separate classical,
coherent, and thermal light.  Glauber received the 2005 Nobel Prize in
Physics for this framework.

The decisive evidence for field quantisation is photon
\emph{antibunching}: the observation that successive photons from a
single quantum emitter are \emph{less} likely to arrive together than
Poisson statistics predict.  \citet{Kimble1977} first observed
sub-Poissonian photon-number statistics in the resonance fluorescence of
individual sodium atoms.  Antibunching is impossible for any classical stochastic field: such a field,
however modulated or noisy, yields a photocount distribution at least as broad
as Poisson, so sub-Poissonian statistics have no classical representation.
The observation therefore refutes the \emph{neoclassical} radiation theories
that sought to dispense with field quantisation altogether, and establishes
that the radiation field is quantised.

The result does not refute \citeauthor{Lamb1969}.  He argued that the
photoelectric effect by itself does not compel field quantisation, and
antibunching leaves that claim standing.  The measurement of
\citet{Kimble1977} required correction for fluctuations in the number of atoms
in the observation volume, which by themselves produce super-Poissonian
counts.  Cleaner demonstrations with single trapped ions followed in the
1980s.

Einstein's 1905 paper introduced
the photon as a \emph{heuristic}---his own word in the title---and
Millikan confirmed its equation to better than $1\%$ in 1916.  Definitive experimental
evidence that the radiation field is itself quantised arrived only in 1977,
by a route---photon-number statistics in resonance fluorescence---that
Einstein's equation neither anticipated nor directly implied.  The
regression confirmed the equation; quantum optics experiments seventy
years later confirmed the physical interpretation.

\subsection{Implications for the Machine Learning Reinterpretation}
\label{app:ml_implications}

The modern results above bear on Section~\ref{sec:ml}.

\paragraph{Depth of the hypothesis class.}
Einstein chose hypothesis class~\eqref{eq:hyp_class} by a heuristic
argument whose full justification---QED, Fermi's golden rule, the
Fermi-Dirac distribution of electrons in a metal---emerged over the
following three decades.  The class proved correct for reasons deeper
than those available in 1905, and it extends over a five-order-of-magnitude
range of photon energies, from visible UV to hard X-rays, without
modification.

\paragraph{What the regression tests.}
Millikan's regression tests the functional form $V_0 \propto \nu$ and
estimates the slope $h/e$.  Both are consistent with the quantised-field and
with the semiclassical model of radiation.  The regression identifies the
best-fitting model within a hypothesis class; selecting between the two
classes requires a different observation, photon-number variance, and a
different experimental design.

\paragraph{Theory specificity in experimental design.}
The antibunching experiments of \citet{Kimble1977} were designed
specifically to test the quantised-field hypothesis; the relevant
observable (photon-number variance) follows directly from
Glauber's coherence theory.  Without the quantised-field hypothesis and
its prediction of sub-Poissonian statistics, there was no reason to
measure photon-number correlations at all.  Theory specifies the hypothesis
class and, from it, the decisive test.  Each theoretical refinement between
Millikan (1916) and \citeauthor{Kimble1977} (1977) generated a new and more
discriminating experiment.  Theory forged the links in that chain of
hypothetico-deductive loops (Table~\ref{tab:loop}).

\bigskip
\bibliographystyle{plainnat}
\bibliography{einstein_photoelectric_ML}

\end{document}
